% Options for packages loaded elsewhere
\PassOptionsToPackage{unicode}{hyperref}
\PassOptionsToPackage{hyphens}{url}
\documentclass[
  ignorenonframetext,
]{beamer}
\newif\ifbibliography
\usepackage{pgfpages}
\setbeamertemplate{caption}[numbered]
\setbeamertemplate{caption label separator}{: }
\setbeamercolor{caption name}{fg=normal text.fg}
\beamertemplatenavigationsymbolsempty
% remove section numbering
\setbeamertemplate{part page}{
  \centering
  \begin{beamercolorbox}[sep=16pt,center]{part title}
    \usebeamerfont{part title}\insertpart\par
  \end{beamercolorbox}
}
\setbeamertemplate{section page}{
  \centering
  \begin{beamercolorbox}[sep=12pt,center]{section title}
    \usebeamerfont{section title}\insertsection\par
  \end{beamercolorbox}
}
\setbeamertemplate{subsection page}{
  \centering
  \begin{beamercolorbox}[sep=8pt,center]{subsection title}
    \usebeamerfont{subsection title}\insertsubsection\par
  \end{beamercolorbox}
}
% Prevent slide breaks in the middle of a paragraph
\widowpenalties 1 10000
\raggedbottom
\AtBeginPart{
  \frame{\partpage}
}
\AtBeginSection{
  \ifbibliography
  \else
    \frame{\sectionpage}
  \fi
}
\AtBeginSubsection{
  \frame{\subsectionpage}
}
\usepackage{iftex}
\ifPDFTeX
  \usepackage[T1]{fontenc}
  \usepackage[utf8]{inputenc}
  \usepackage{textcomp} % provide euro and other symbols
\else % if luatex or xetex
  \usepackage{unicode-math} % this also loads fontspec
  \defaultfontfeatures{Scale=MatchLowercase}
  \defaultfontfeatures[\rmfamily]{Ligatures=TeX,Scale=1}
\fi
\usepackage{lmodern}
\ifPDFTeX\else
  % xetex/luatex font selection
\fi
% Use upquote if available, for straight quotes in verbatim environments
\IfFileExists{upquote.sty}{\usepackage{upquote}}{}
\IfFileExists{microtype.sty}{% use microtype if available
  \usepackage[]{microtype}
  \UseMicrotypeSet[protrusion]{basicmath} % disable protrusion for tt fonts
}{}
\makeatletter
\@ifundefined{KOMAClassName}{% if non-KOMA class
  \IfFileExists{parskip.sty}{%
    \usepackage{parskip}
  }{% else
    \setlength{\parindent}{0pt}
    \setlength{\parskip}{6pt plus 2pt minus 1pt}}
}{% if KOMA class
  \KOMAoptions{parskip=half}}
\makeatother
\setlength{\emergencystretch}{3em} % prevent overfull lines
\providecommand{\tightlist}{%
  \setlength{\itemsep}{0pt}\setlength{\parskip}{0pt}}
\usepackage{bookmark}
\IfFileExists{xurl.sty}{\usepackage{xurl}}{} % add URL line breaks if available
\urlstyle{same}
\hypersetup{
  hidelinks,
  pdfcreator={LaTeX via pandoc}}

\author{\texorpdfstring{}{}}
\date{}

\begin{document}

\begin{frame}
\vspace*{2cm}

\begin{center}
\begin{minipage}{0.7\textwidth}
\centering
\Large\itshape
``In theory, there is no difference between\\[0.3em]
theory and practice. In practice, there is.''
\vspace{1em}

\normalsize\upshape
--- Yogi Berra (attributed)
\end{minipage}
\end{center}

\vspace{2cm}
\end{frame}

\begin{frame}{Abstract}
\protect\phantomsection\label{abstract}
Multiagent AI systems---autonomous coding assistants, research
pipelines, financial decision engines---have moved from prototype to
production in under two years. With them comes a new class of security
concern: attacks that target not data or infrastructure but the
\emph{reasoning processes} of AI agents. Prompt injections that
propagate through delegation chains, trust relationships that launder
adversarial influence, and coordination mechanisms vulnerable to
strategic manipulation all represent cognitive attack surfaces absent
from traditional security models.

The \textbf{Cognitive Integrity Framework (CIF)}, presented across two
companion papers, offers the first comprehensive formal and
computational treatment of this problem. Part 1 establishes mathematical
foundations: a trust calculus with provably bounded delegation, defense
composition algebras with multiplicative detection guarantees, and
information-theoretic limits on attack stealth. Part 2 provides
computational validation: eight implemented defense modules (1,557
passing tests), a 950-attack corpus spanning four threat categories, and
parametric architecture-aware simulation across six production
multiagent topologies.

This paper is a qualitative review and practitioner's guide to the CIF
series. We synthesize the key insights from both papers into accessible
language, contextualize the formal results within the current deployment
landscape, assess what the research has established and where gaps
remain, and distill practical recommendations for teams building and
operating multiagent AI systems. No formal prerequisites are assumed;
readers seeking mathematical detail are referred to Parts 1 and 2.

\begin{block}{Paper Series}
\protect\phantomsection\label{paper-series}
\textbf{DOI}: 10.5281/zenodo.18364130

This is Part 3 of the \emph{Cognitive Security for Multiagent Operators}
series:

\begin{itemize}
\tightlist
\item
  \textbf{Part 1} (DOI: 10.5281/zenodo.18364119): Formal foundations and
  theoretical analysis
\item
  \textbf{Part 2} (DOI: 10.5281/zenodo.18364128): Computational
  validation and implementation
\item
  \textbf{Part 3} (this paper): Qualitative review and practitioner's
  synthesis
\end{itemize}
\end{block}
\end{frame}

\end{document}
